% Subsection: Examples of Seed Families  (notes stub; content authored later)
\subsection{Examples of Seed Families}

\begin{remark*}[Exposition]
Most useful set systems are not specified by listing every admitted subset.
Instead, one begins with a small family whose members are easy to describe,
then enlarges it by enforcing chosen closure rules.  The initial family is the
seed data.  Later sections make the generated enlargement precise.
\end{remark*}

\begin{definition}[Seed Family]\label{def:seed-family}
Let \(X\) be a set.  A \emph{seed family on \(X\)} is a family
\(\mathcal{G}\subseteq\mathcal{P}(X)\) chosen as initial data for constructing
or generating a larger set system on \(X\).
\end{definition}

\begin{remark*}[Standard quantified statement]
\[
  \mathcal{G}\text{ is a seed family on }X
  \quad\Longrightarrow\quad
  \mathcal{G}\subseteq\mathcal{P}(X).
\]
\end{remark*}

\begin{remark*}[Interpretation]
The word ``seed'' does not add a new set-theoretic condition beyond being a
family of subsets.  It records the intended use of the family: the members of
\(\mathcal{G}\) are the starting pieces from which a larger admissible system
will be generated.
\end{remark*}

\begin{remark*}[Exposition]
The same pattern will recur throughout the book.  A family of basic sets is
chosen first; then a structural operation closes that family into the kind of
system needed for the theory.  Metric-space chapters will use balls as basic
sets.  Topological chapters will study systems of open sets.  Measure-theory
chapters will use sigma-algebras to specify the measurable subsets.
\end{remark*}

\begin{remark*}[Examples]
Typical seed families include intervals in \(\mathbb{R}\), rectangles in
\(\mathbb{R}^{n}\), balls in a metric setting, and any named collection of
subsets whose closure under later operations is intended to produce a larger
system.  At this stage these are only families of subsets.  Their later role
depends on which closure rule is imposed.
\end{remark*}

\begin{dependencies}
\begin{itemize}
  \item \hyperref[def:family-of-subsets]{Family of Subsets}
\end{itemize}
\end{dependencies}
